% TEMPLATE-MILC-v3.tex - plantilla maestra para todas las guias MILC v3
% Ver scripts/generadores/build-guia-milc.py para la lista completa de
% claves esperadas. El script valida que no queden placeholders sin
% reemplazar antes de escribir el archivo final.

\documentclass[11pt,letterpaper]{article}
\usepackage[margin=1.75cm]{geometry}
\usepackage[table]{xcolor}
\usepackage{fontspec}
\usepackage[spanish]{babel}
\usepackage{tikz}
\usepackage[most]{tcolorbox}
\usepackage{tabularx}
\usepackage{array}
\usepackage{fancyhdr}
\usepackage{titlesec}
\usepackage{ragged2e}
\usepackage{enumitem}
\usepackage{microtype}

\setmainfont{Helvetica Neue}
\setsansfont{Helvetica Neue}
\setlength{\parindent}{0pt}
\setlength{\parskip}{6pt}
\hyphenpenalty=200
\exhyphenpenalty=200
\pretolerance=400
\tolerance=2000
\setlength{\emergencystretch}{1em}
\renewcommand{\arraystretch}{1.25}
\setlist{leftmargin=5mm,itemsep=1mm,topsep=1mm}
\newcolumntype{Y}{>{\RaggedRight\arraybackslash}X}

\definecolor{milcMagenta}{HTML}{E6007E}
\definecolor{milcVino}{HTML}{5A0038}
\definecolor{milcTurquesa}{HTML}{0093A5}
\definecolor{milcVerde}{HTML}{58A923}
\definecolor{milcMostaza}{HTML}{E5B400}
\definecolor{milcOcre}{HTML}{B86600}
\definecolor{milcNegro}{HTML}{241816}
\definecolor{milcGris}{HTML}{F7F7F4}
\definecolor{milcLinea}{HTML}{E8E2DD}
\definecolor{milcGradePrimary}{HTML}{0066FF}
\definecolor{milcGradeDark}{HTML}{003D99}
\definecolor{milcGradeSoft}{HTML}{D6E8FF}
\definecolor{milcGradeAccent}{HTML}{CFE7FF}
\definecolor{milcGradeText}{HTML}{FFFFFF}
\color{milcNegro}

\pagestyle{fancy}
\fancyhf{}
\lhead{\small Tecnología e Informática · Grado 11}
\rhead{\small Guía 12 · MILC v3}
\cfoot{\small\thepage}
\renewcommand{\headrulewidth}{0pt}

\newtcolorbox{titlebox}[1]{
  enhanced,colback=#1,colframe=#1,arc=7mm,boxrule=0pt,
  left=7mm,right=7mm,top=3mm,bottom=3mm,
  before skip=8pt,after skip=8pt,
  fontupper=\sffamily\bfseries\Large\color{white}
}

\newtcolorbox{softbox}[3]{
  enhanced,colback=#2,colframe=#1,borderline west={4pt}{0pt}{#1},
  arc=2mm,boxrule=.4pt,left=4mm,right=4mm,top=3mm,bottom=3mm,
  before skip=6pt,after skip=8pt,title={#3},coltitle=white,
  colbacktitle=#1,fonttitle=\sffamily\bfseries,fontupper=\RaggedRight,
  attach boxed title to top left={xshift=3mm,yshift=-2mm},
  boxed title style={arc=2mm, boxrule=0pt}
}

\newtcolorbox{drawbox}[2]{
  enhanced,colback=white,colframe=#1,arc=4mm,boxrule=1pt,
  left=5mm,right=5mm,top=4mm,bottom=4mm,before skip=8pt,after skip=8pt,
  title={#2},coltitle=white,colbacktitle=#1,fonttitle=\sffamily\bfseries,
  fontupper=\RaggedRight,
  attach boxed title to top left={xshift=4mm,yshift=-2mm},
  boxed title style={arc=3mm, boxrule=0pt}
}

\newtcolorbox{infoband}[1]{
  enhanced,colback=#1!8,colframe=#1!50,arc=2mm,boxrule=.5pt,
  left=4mm,right=4mm,top=2mm,bottom=2mm,before skip=4pt,after skip=4pt,
  fontupper=\small\RaggedRight
}

% Puente narrativo entre secciones (contrato editorial Conéctate).
% Da continuidad: "Acabas de X. Ahora vas a Y porque Z."
% Visualmente: banda izquierda en color del grado, texto en itálica pequeña.
\newtcolorbox{puentebox}{
  enhanced,colback=milcGradeSoft,colframe=milcGradeDark,
  borderline west={3pt}{0pt}{milcGradeDark},
  arc=0mm,boxrule=0pt,left=5mm,right=4mm,top=2.5mm,bottom=2.5mm,
  before skip=4pt,after skip=8pt,
  fontupper=\itshape\small\color{milcNegro}\RaggedRight
}
\newcommand{\puente}[1]{%
  \begin{puentebox}%
    \textcolor{milcGradeDark}{\textbf{\textrightarrow}}\hspace{2mm}#1%
  \end{puentebox}%
}

\newcommand{\linead}[1]{\noindent\color{milcLinea}\rule{#1}{0.5pt}\color{milcNegro}}
\newcommand{\respuesta}[1]{\par\vspace{2mm}\foreach \n in {1,...,#1}{\noindent\color{milcLinea}\rule{\linewidth}{0.5pt}\par\vspace{5mm}}\color{milcNegro}}
\newcommand{\checkbox}{\textcolor{milcGradeDark}{\fbox{\phantom{x}}}\hspace{5pt}}

\newcommand{\phasecard}[4]{
\begin{tcolorbox}[enhanced,colback=#3,colframe=#2,arc=3mm,boxrule=.5pt,height=3.05cm,valign=center,halign=center,left=1.5mm,right=1.5mm,top=2mm,bottom=2mm]
{\sffamily\bfseries\fontsize{22}{24}\selectfont\textcolor{#2}{#1}}\par
{\sffamily\bfseries\small #4}
\end{tcolorbox}
}

\begin{document}
\thispagestyle{empty}
\begin{tikzpicture}[remember picture,overlay]
  \fill[white] (current page.south west) rectangle (current page.north east);
  \path[fill=milcGradePrimary,rounded corners=16mm]
    ([xshift=.55cm,yshift=-.55cm]current page.north west) rectangle
    ([xshift=-.55cm,yshift=3.65cm]current page.south east);

  \node[anchor=north west,text=milcGradeText] at ([xshift=2.0cm,yshift=-1.85cm]current page.north west)
    {\sffamily\bfseries\fontsize{14}{16}\selectfont GRADO};
  \node[anchor=north west,text=milcGradeText,opacity=.82] at ([xshift=2.0cm,yshift=-2.75cm]current page.north west)
    {\sffamily\bfseries\fontsize{9}{11}\selectfont SERIE GUÍAS · TECNOLOGÍA E INFORMÁTICA};

  \begin{scope}[shift={([xshift=-3.05cm,yshift=-2.55cm]current page.north east)}]
    \fill[white,opacity=.80] (-.62,-.52) rectangle (.62,.52);
    \draw[milcGradeText,line width=1pt] (-.62,-.52) rectangle (.62,.52);
    \fill[milcGradeDark] (-.42,-.38) rectangle (-.22,.06);
    \fill[milcGradeAccent] (-.10,-.38) rectangle (.10,.24);
    \fill[milcGradePrimary!60!black] (.22,-.38) rectangle (.42,.38);
  \end{scope}

  \node[anchor=west,text=milcGradeText] at ([xshift=1.9cm,yshift=-12.85cm]current page.north west)
    {\sffamily\bfseries\fontsize{124}{124}\selectfont 11};
  \draw[milcGradeText,line width=1.7pt]
    ([xshift=4.52cm,yshift=-11.68cm]current page.north west) circle (.13cm);

  \node[anchor=west,text=milcGradeText,text width=8.7cm,align=left] at ([xshift=10.35cm,yshift=-11.6cm]current page.north west)
    {
     {\sffamily\bfseries\fontsize{19}{22}\selectfont Diagramas de flujo ---\\lenguaje\\BPMN básico}\\[4mm]
     {\sffamily\bfseries\fontsize{23}{26}\selectfont Guía 12-11-TIC}\\[2mm]
     {\sffamily\fontsize{13}{16}\selectfont Período 2 · Automatización y procesos}\\[5mm]
     {\sffamily\bfseries\fontsize{11}{13}\selectfont Institución Educativa Sor María Juliana}\\[1mm]
     {\sffamily\fontsize{10}{12}\selectfont Autor: Dr. Álvaro Cárdenas Orozco}
    };

  \path[fill=milcGradeSoft,rounded corners=7mm]
    ([xshift=1.35cm,yshift=.85cm]current page.south west) rectangle
    ([xshift=-1.35cm,yshift=4.25cm]current page.south east);
  \draw[milcGradeDark,line width=.65pt,rounded corners=7mm]
    ([xshift=1.35cm,yshift=.85cm]current page.south west) rectangle
    ([xshift=-1.35cm,yshift=4.25cm]current page.south east);
  \node[anchor=center,text=milcNegro,text width=17.0cm,align=center] at ([yshift=2.55cm]current page.south)
    {\sffamily\fontsize{10}{12}\selectfont Metodología MILC · Apertura ancestral · Pensamiento computacional · Triángulo Dussel-Estoicismo-Floridi\\
    \textcolor{milcGradeDark}{\bfseries Producto:} Diagrama BPMN del proceso mapeado en sesión 1, hecho en draw.io o Lucidchart, usando correctamente los 5 símbolos básicos (inicio/fin, tarea, decisión, flecha, swimlane) y exportado en PDF para entrega};
\end{tikzpicture}
\mbox{}
\newpage

\section*{Apertura: Diagramas de flujo --- lenguaje BPMN básico}

\begin{softbox}{milcGradeDark}{milcGradeSoft}{Saber ancestral}
Las \textbf{recetas familiares} del Valle son diagramas de flujo escritos sin saberlo: \emph{``primero el sofrito, luego el caldo, después el grano, al final el cilantro''}. La \textbf{cocinera de la fonda} sabe que cambiar el orden daña el sancocho; el \textbf{maestro carpintero} sabe que lijar antes de cortar arruina la madera; el \textbf{tejedor Wayuu} sabe que cada rombo del chinchorro se hace en una secuencia precisa, no negociable. Las recetas, los manuales del oficio y los pasos heredados son \emph{flujos formalizados por la experiencia colectiva}. BPMN no inventó los diagramas de flujo: les puso un alfabeto común para que un ingeniero en Bogotá entienda el flujo de una fábrica en Cartago sin necesidad de visitarla.
\end{softbox}

\begin{softbox}{milcTurquesa}{milcGris}{Saber contemporáneo}
\textbf{BPMN} (\emph{Business Process Model and Notation}) es un estándar internacional para diagramar procesos de negocio. Como las partituras musicales o los planos arquitectónicos, BPMN tiene un \emph{vocabulario cerrado} que cualquier profesional del mundo lee igual. Los 5 símbolos básicos del nivel introductorio son: (1) \textbf{círculo}: inicio (verde) o fin (rojo); (2) \textbf{rectángulo}: tarea o actividad; (3) \textbf{rombo}: decisión (\emph{``¿sí o no?''}); (4) \textbf{flecha}: flujo de secuencia; (5) \textbf{swimlane}: carril que agrupa por responsable. Dominar estos 5 símbolos te permite leer y producir el 80\% de los diagramas profesionales.
\end{softbox}

\begin{softbox}{milcMostaza}{milcGris}{Pregunta-puente}
\textit{¿Cómo sabía la cocinera, sin haberlo escrito, qué pasos del sancocho admiten desorden y cuáles no? ¿Y qué pierde un emprendedor novato cuando inventa su propio \emph{``lenguaje gráfico''} en lugar de aprender el estándar que todos los profesionales del mundo ya leen?}
\end{softbox}

\begin{softbox}{milcVerde}{milcGris}{Saber y saber hacer}
Al terminar podrás: (1) \textbf{identificar} los 5 símbolos básicos de BPMN y qué representa cada uno; (2) \textbf{explicar} con tus palabras la diferencia entre tarea, decisión, evento y swimlane; (3) \textbf{aplicar} BPMN al proceso que mapeaste en sesión 1, convirtiendo tu boceto en un diagrama profesional; (4) \textbf{evaluar} si tu diagrama es legible para alguien que NO conoce tu proceso (criterio de profesionalidad).
\end{softbox}

\small
\begin{tabularx}{\textwidth}{>{\bfseries}p{3.0cm}Y}
\rowcolor{milcGradePrimary!12}Autor & Dr. Álvaro Cárdenas Orozco\\
DBA & Diseño y mejora de procesos digitales con criterio ético (MEN, T\&I)\\
Referentes & Pensamiento computacional · Lean / Kaizen · Floridi (ética informacional)\\
Producto final & Diagrama BPMN del proceso mapeado en sesión 1, hecho en draw.io o Lucidchart, usando correctamente los 5 símbolos básicos (inicio/fin, tarea, decisión, flecha, swimlane) y exportado en PDF para entrega\\
\end{tabularx}
\normalsize

\puente{Antes de empezar, mira el viaje completo. En la sesión 1 mapeaste a mano alzada un proceso real con sus cuellos de botella. Hoy le pones \emph{lenguaje profesional}: el mismo proceso, pero diagramado en BPMN para que un consultor, un ingeniero o un emprendedor lo lea sin verte explicar.}

\begin{titlebox}{milcGradeDark}
Ruta MILC: así vamos a aprender
\end{titlebox}
\begin{tabularx}{\textwidth}{>{\centering\arraybackslash}X>{\centering\arraybackslash}X>{\centering\arraybackslash}X>{\centering\arraybackslash}X}
\phasecard{1}{milcVerde}{milcGris}{Escuta\\[-1mm]\small Observo, escucho y pregunto.} &
\phasecard{2}{milcTurquesa}{milcGris}{Sistematización\\[-1mm]\small Pensamiento computacional.} &
\phasecard{3}{milcMagenta}{milcGris}{Praxis\\[-1mm]\small Construyo y aplico.} &
\phasecard{4}{milcVino}{milcGris}{Evaluación\\[-1mm]\small Reflexión liberadora.}
\end{tabularx}


\puente{Antes de teorizar, vas a leer 3 diagramas BPMN reales sin instrucción previa y a inferir qué hace cada uno. Esa lectura sin guía mide qué tan intuitivo es el lenguaje (mucho más de lo que parece a primera vista).}

\begin{titlebox}{milcVerde}
Fase 1 · Escuta
\end{titlebox}

\begin{softbox}{milcVerde}{milcGris}{Escena para escuchar}
\textbf{Actividad 1 · IDENTIFICA --- Lee 3 BPMN sin guía} (15 min · individual). El docente proyecta 3 diagramas BPMN reales (ej. \emph{registro de un nuevo cliente en banco}, \emph{pedido en una pizzería}, \emph{trámite de matrícula escolar}). Sin saber qué símbolo significa qué, intenta describir con tus palabras qué hace cada diagrama: ¿dónde empieza?, ¿qué decisiones aparecen?, ¿quiénes están involucrados?, ¿dónde termina? Anota lo que entiendes y lo que te confunde.
\end{softbox}

\checkbox Describí en mis palabras qué hace cada uno de los 3 diagramas.\\
\checkbox Identifiqué inicio, fin y al menos 1 decisión en cada diagrama.\\
\checkbox Anoté qué símbolo o flecha me resultó confusa para preguntar después.

\begin{infoband}{milcVerde}
\textbf{Lo que va al cuaderno (Actividad 1 · IDENTIFICA).} Título: «Actividad 1 --- Lee 3 BPMN sin guía». \textbf{Formato:} tabla 3 filas (un diagrama por fila) × 3 columnas (Qué hace el proceso~|~Decisiones que detecté~|~Lo que me confunde). \textbf{Sabes que terminaste cuando} reconoces que BPMN es más intuitivo de lo que parece y puedes señalar dónde te trabaste.
\end{infoband}

\puente{Ya descifraste 3 diagramas. Ahora vamos a entender la \textbf{anatomía} de los 5 símbolos básicos y las reglas mínimas que distinguen un BPMN profesional de un dibujito desordenado: un solo inicio, un solo fin (idealmente), flechas que no se cruzan si se puede evitar, decisiones que siempre dan dos salidas etiquetadas.}

\begin{titlebox}{milcTurquesa}
Fase 2 · Sistematización con pensamiento computacional
\end{titlebox}

\begin{softbox}{milcTurquesa}{milcGris}{Cuatro pilares aplicados al tema}
Los 5 símbolos básicos de BPMN se aprenden en 15 minutos pero el rigor de su uso es lo que distingue un diagrama profesional. Vamos a descomponerlos con los pilares del pensamiento computacional.
\end{softbox}

\small
\begin{tabularx}{\textwidth}{>{\bfseries\RaggedRight\arraybackslash}p{3.6cm}Y}
\rowcolor{milcTurquesa!90}\color{white}Pilar & \color{white}\bfseries Aplicación al tema\\
1. Descomposición & Los 5 símbolos se descomponen así: (1) \textbf{Eventos} (círculos): inicio (verde, fino), fin (rojo, grueso), evento intermedio (doble línea); marcan el comienzo y el cierre de un flujo. (2) \textbf{Tareas} (rectángulos redondeados): cada acción concreta del proceso (\emph{``recibir pedido''}, \emph{``verificar stock''}). (3) \textbf{Compuertas} (rombos): decisiones; deben tener mínimo 2 salidas etiquetadas con condición (\emph{``sí / no''}, \emph{``pago aprobado / rechazado''}). (4) \textbf{Flujos de secuencia} (flechas continuas): qué pasa después de qué. (5) \textbf{Swimlanes} (carriles horizontales): agrupan tareas por responsable (\emph{cliente, vendedor, sistema}).\\
2. Reconocimiento de patrones & Todo BPMN sigue el patrón \textbf{``un solo inicio, un solo fin (idealmente)''}: si tu diagrama tiene 3 puntos de inicio, probablemente son 3 procesos distintos mezclados. Si tiene 4 fines, alguno es solo un cierre parcial. La regla profesional: \emph{procesos limpios = pocos eventos extremos}. Los expertos a veces rompen esta regla con razón, pero al principiante le sirve como disciplina.\\
3. Abstracción & Lo esencial cabe en una pregunta: \emph{¿alguien que no conoce mi proceso puede seguir las flechas y entender el flujo completo?}. Si la respuesta es no, el diagrama no es profesional aunque sea bonito. La prueba final es la legibilidad ciega: \emph{darle el diagrama a un compañero sin explicación}.\\
4. Algoritmo conceptual & Algoritmo para construir un BPMN: (1) lista responsables (van en swimlanes); (2) marca inicio (círculo verde); (3) traduce cada paso a tarea (rectángulo); (4) cada \emph{``si X, hacemos A, si Y, hacemos B''} a compuerta (rombo); (5) conecta con flechas; (6) marca fin (círculo rojo); (7) revisa que no haya flechas sueltas ni cruces evitables.\\
\end{tabularx}
\normalsize

\begin{drawbox}{milcTurquesa}{Anatomía de un BPMN profesional}
\textbf{Título:} nombre del proceso. \\[1mm] \textbf{Swimlanes:} responsables como carriles horizontales. \\[1mm] \textbf{Inicio:} círculo verde único. \\[1mm] \textbf{Tareas:} rectángulos con verbo de acción y objeto (\emph{``verificar stock''}, no \emph{``stock''}). \\[1mm] \textbf{Compuertas:} rombos con condición + 2 salidas etiquetadas. \\[1mm] \textbf{Fin:} círculo rojo (idealmente uno). \\[1mm] \textbf{Flechas:} continuas, sin cruces evitables.
\end{drawbox}

\begin{softbox}{milcMostaza}{milcGris}{Errores comunes y su corrección}
(1) Rombos con una sola salida: no son decisión, son tareas mal dibujadas. (2) Tareas con sustantivos en lugar de verbos: \emph{``factura''} no es tarea; \emph{``emitir factura''} sí. (3) Flechas que no etiquetan las decisiones: el lector no sabe cuál camino es cuál. (4) Múltiples inicios sin justificación: confunde al lector. (5) Olvidar swimlanes: si hay más de 1 responsable, los carriles son obligatorios para legibilidad.
\end{softbox}

\begin{infoband}{milcTurquesa}
\textbf{Lo que va al cuaderno (Actividad 2 · EXPLICA).} Título: «Actividad 2 --- Anatomía BPMN». \textbf{Formato:} ficha con los 5 símbolos dibujados + qué representa cada uno + 1 error frecuente por símbolo. \textbf{Sabes que terminaste cuando} podrías corregir un BPMN ajeno señalando con vocabulario técnico qué está mal.
\end{infoband}

\puente{Tienes el método. Toca aplicarlo: abrir draw.io o Lucidchart, recrear el proceso de sesión 1 con los 5 símbolos, organizar por swimlanes según responsables, y exportar en PDF.}

\begin{titlebox}{milcMagenta}
Fase 3 · Praxis con pensamiento computacional
\end{titlebox}

\begin{softbox}{milcMagenta}{milcGris}{Construyendo el producto con los cuatro pilares}
\textbf{Actividad 3 · APLICA --- Tu proceso en BPMN profesional} (30 min · individual). Hora de aplicar. Abre draw.io (gratuito, sin cuenta) o Lucidchart, recrea el proceso de sesión 1 con los 5 símbolos básicos, organiza por swimlanes según responsables, y exporta en PDF.
\end{softbox}

\small
\begin{tabularx}{\textwidth}{>{\bfseries\RaggedRight\arraybackslash}p{3.6cm}Y}
\rowcolor{milcMagenta!90}\color{white}Pilar & \color{white}\bfseries Aplicación al producto\\
Descomposición & Tu diagrama se descompone en 5 piezas: (1) título del proceso al inicio del lienzo; (2) swimlanes con cada responsable; (3) inicio + tareas + compuertas + fin correctamente ubicados en su carril; (4) flechas etiquetadas en las decisiones; (5) exportación en PDF lista para entrega.\\
Patrones & Sigue el patrón \textbf{``flujo de izquierda a derecha, sin retrocesos visuales innecesarios''}: el lector recorre el diagrama como un texto. Si una flecha vuelve atrás, debe ser un bucle real (\emph{``rechazado $\to$ corregir''}), no un capricho de diseño. Cuanto más lineal el flujo, más legible.\\
Abstracción & Cada compuerta debe etiquetar sus salidas: \emph{sí/no}, \emph{aprobado/rechazado}, \emph{cliente nuevo/cliente recurrente}. Si una compuerta tiene 3 salidas, las 3 deben tener etiqueta. El lector profesional debe poder seguir cada rama sin adivinar.\\
Algoritmo de armado & Algoritmo: (1) abre draw.io; (2) crea swimlanes; (3) coloca inicio en el carril correspondiente; (4) arrastra tareas, compuertas, flechas; (5) etiqueta cada decisión; (6) ubica el fin; (7) revisa cruces y limpieza visual; (8) exporta a PDF.\\
\end{tabularx}
\normalsize

\begin{drawbox}{milcMagenta}{Checklist antes de entregar el BPMN}
\checkbox Título del proceso visible. \\[1mm] \checkbox Swimlanes por responsable si hay más de uno. \\[1mm] \checkbox Inicio único (círculo verde) y fin único (círculo rojo). \\[1mm] \checkbox Tareas con verbo + objeto. \\[1mm] \checkbox Compuertas con condición y salidas etiquetadas. \\[1mm] \checkbox Sin flechas sueltas, sin cruces innecesarios. \\[1mm] \checkbox Exportado como PDF legible.
\end{drawbox}

\begin{softbox}{milcMagenta}{milcGris}{Plantilla de trabajo}
\textbf{Proceso.} \textit{Nombre + responsable general.} \\[1mm] \textbf{Swimlanes.} \textit{Cliente / Vendedor / Sistema / etc.} \\[1mm] \textbf{Inicio.} \textit{Evento que dispara el proceso.} \\[1mm] \textbf{Tareas principales.} \textit{5-10 tareas con verbo + objeto.} \\[1mm] \textbf{Compuertas.} \textit{2-3 decisiones con salidas etiquetadas.} \\[1mm] \textbf{Fin.} \textit{Estado final del proceso.}
\end{softbox}

\begin{infoband}{milcMagenta}
\textbf{Lo que va al cuaderno (Actividad 3 · APLICA).} Título: «Actividad 3 --- BPMN de mi proceso». \textbf{Formato:} pantallazo del diagrama pegado o impreso + checklist completa marcada + 1 línea con qué fue lo más difícil. \textbf{Extensión:} 1 página de cuaderno con el diagrama. \textbf{Sabes que terminaste cuando} un compañero que no vio el proceso original puede explicarte cómo funciona leyendo solo tu diagrama.
\end{infoband}

\puente{Lo que sigue resume \emph{qué} entregas hoy: 1 diagrama BPMN en PDF + 1 archivo editable. El criterio principal: que un profesional que NO conoce tu proceso pueda \textbf{entenderlo} solo con el diagrama, sin tu explicación.}

\begin{titlebox}{milcMostaza}
Producto de la sesión
\end{titlebox}
\textbf{1 diagrama BPMN profesional en draw.io/Lucidchart, exportado a PDF}.
\begin{softbox}{milcMostaza}{milcGris}{Criterios de calidad}
(1) Título y swimlanes presentes. (2) 5 símbolos básicos usados con corrección. (3) Inicio y fin únicos. (4) Compuertas con salidas etiquetadas. (5) Exportado a PDF, archivo editable conservado.
\end{softbox}

\puente{Antes de cerrar, mira el diagrama desde las cinco dimensiones humanas. Un BPMN bien hecho es respeto al lector profesional; uno mal hecho es desperdicio de tiempo colectivo.}

\begin{titlebox}{milcVino}
Fase 4 · Evaluación liberadora — cinco dimensiones
\end{titlebox}

\begin{softbox}{milcVino}{milcGris}{Sentido de la evaluación}
Evaluar no es solo revisar una nota. Es reconocer qué aprendí, qué cambió en mí, qué emociones manejé, cómo me relaciono con la ciudad y la comunidad, y qué le dejaré a quien viene detrás.
\end{softbox}

\textbf{Mi reflexión liberadora en cinco dimensiones.} Completa cada frase iniciada:

\small
\begin{tabularx}{\textwidth}{>{\bfseries\RaggedRight\arraybackslash}p{3.4cm}Y>{\RaggedRight\arraybackslash}p{4.6cm}}
\rowcolor{milcVino}\color{white}Dimensión & \color{white}\bfseries Mi respuesta (frame starter) & \color{white}\bfseries Algo que descubrí\\
1. Personal & Lo que más cambió en mí fue \linead{6.5cm} & \linead{4.2cm}\\
2. Emocional & La emoción más fuerte fue \linead{2.5cm} y la regulé \linead{3cm} & \linead{4.2cm}\\
3. Ciudadana & En democracia esto importa porque \linead{6cm} & \linead{4.2cm}\\
4. Local (Cartago) & En mi barrio sería útil para \linead{6.5cm} & \linead{4.2cm}\\
5. Intergeneracional & A mi abuela/abuelo le diría \linead{6.5cm} & \linead{4.2cm}\\
\end{tabularx}
\normalsize

\begin{infoband}{milcVino}
\textbf{Para la próxima sustentación.} Vuelve a esta tabla en seis meses. Verás que las respuestas cambian — no porque lo aprendido cambie, sino porque tú cambias. La evaluación liberadora es una conversación contigo a través del tiempo.
\end{infoband}


\puente{Y para terminar, tres voces sobre el lenguaje compartido. Dussel sobre los lenguajes técnicos que excluyen, Epicteto sobre la disciplina del estándar, Floridi sobre BPMN como infraestructura de la economía digital.}

\begin{titlebox}{milcGradeDark}
Triángulo de pensamiento · cierre filosófico
\end{titlebox}

\begin{softbox}{milcVino}{milcGris}{Dussel · lente del nosotros}
\textit{``Todo lenguaje técnico decide quién entra a la conversación y quién queda fuera por desconocer el alfabeto.''}

BPMN abre puertas profesionales, pero también las cierra a quienes no han tenido acceso a la formación. La phronesis del que aprende el lenguaje técnico exige: (1) aprenderlo bien para no quedar excluido; (2) traducirlo a quien no lo conoce para no excluir. Un BPMN bonito que el responsable real del proceso no entiende es un BPMN mudo. Saber BPMN obliga a saber explicarlo en castellano sencillo cuando el contexto lo pide.

\textbf{¿Puedo explicar mi BPMN a la cocinera real de la fonda o al maestro carpintero real del taller en su lenguaje, o solo lo entienden los \emph{``del oficio digital''}?}
\end{softbox}
\linead{\linewidth}\\[1mm]
\linead{\linewidth}

\begin{softbox}{milcMostaza}{milcGris}{Estoicismo · Epicteto · cuidado interior}
\textit{``Aprender el estándar es aceptar que el oficio es más grande que tu preferencia.''}

Es tentador \emph{``inventar mi propio sistema de diagramación porque me parece más claro''}. La sabiduría estoica recomienda lo contrario: \emph{el estándar es la humildad del oficio profesional}. BPMN no es la única forma posible de diagramar, pero es la que el mundo profesional acordó. Quien se niega a aprenderlo se queda fuera de la conversación; quien lo aprende entra al diálogo de la ingeniería contemporánea.

\textbf{¿Resistí la tentación de inventar mi propio lenguaje y acepté el estándar profesional? ¿Qué libertad gano al hablar la lengua que todos hablan?}
\end{softbox}
\linead{\linewidth}\\[1mm]
\linead{\linewidth}

\begin{softbox}{milcTurquesa}{milcGris}{Floridi · lente de la infoesfera}
\textit{``Los estándares de notación son la infraestructura invisible de toda economía digital colaborativa.''}

Cada vez que un equipo internacional diseña un software, cada vez que un consultor evalúa una empresa, cada vez que un ingeniero entrega un proyecto, BPMN aparece. Es la lengua franca de los procesos. Aprenderlo a los 16 años es alfabetización profesional adelantada: el lenguaje con el que vas a comunicarte en cualquier empleo que involucre procesos, sistemas o transformación digital.

\textbf{¿Cuántas conversaciones profesionales futuras (con un consultor, un ingeniero, un emprendedor) dependerán de que comparta el lenguaje BPMN? ¿Cuánto me cuesta no aprenderlo bien ahora?}
\end{softbox}
\linead{\linewidth}\\[1mm]
\linead{\linewidth}

\begin{titlebox}{milcVino}
Compromiso final
\end{titlebox}

\small
\begin{tabularx}{\textwidth}{>{\RaggedRight\arraybackslash}p{3.0cm}YYY}
\rowcolor{milcVino}\color{white}\bfseries Criterio & \color{white}\bfseries Lo logré & \color{white}\bfseries En proceso & \color{white}\bfseries Necesito apoyo\\
Comprensión del tema & Explico ideas centrales con mis palabras. & Reconozco ideas pero falta precisión. & Necesito repasar conceptos base.\\
Pensamiento computacional & Apliqué los 4 pilares al diseñar mi producto. & Apliqué algunos pilares. & Necesito releer la fase 2.\\
Aplicación y producto & Mi producto cumple los criterios y se puede revisar. & Producto funcional con ajustes pendientes. & Aún en construcción.\\
Reflexión 5 dimensiones & Respondí cada dimensión con honestidad. & Respondí algunas con detalle. & Mis respuestas son muy generales.\\
Triángulo de pensamiento & Conecté las 3 voces con mi trabajo real. & Conecté al menos una. & No relaciono las voces con mi proceso.\\
\end{tabularx}
\normalsize

\textbf{Hoy entendí que:}\\[1mm]
\linead{\linewidth}\\[1mm]
\linead{\linewidth}

\textbf{Mi compromiso para la próxima sesión es:}\\[1mm]
\linead{\linewidth}\\[1mm]
\linead{\linewidth}

\begin{softbox}{milcMostaza}{milcGris}{Cita de despedida · Marco Aurelio}
\textit{``El obstáculo en el camino se vuelve el camino. Lo que estorba al camino se vuelve camino.''}\\[1mm]
Cada error de hoy, bien observado, se convierte en pieza del camino que pisarás mañana.
\end{softbox}

\small
\begin{tabularx}{\textwidth}{>{\bfseries}p{3.6cm}Y}
\rowcolor{milcGradeDark!12}Nombre completo: & \linead{10cm}\\
Fecha: & \linead{4cm} \quad Curso/grupo: \linead{4cm}\\
Mi firma: & \linead{9cm}\\
\end{tabularx}
\normalsize

\begin{infoband}{milcGradeDark}
\textbf{Para la próxima semana.} Toma una receta familiar de tu casa (sancocho, sopa, plato emblema) y diagrámala en BPMN en draw.io. Incluye swimlanes (cocinero principal, ayudante, comensal que pide), decisiones (\emph{``¿hay yuca? sí/no''}) y eventos. La phronesis del lenguaje técnico se entrena traduciéndolo a lo cotidiano.
\end{infoband}

\end{document}
